\section{Pre-registered transfer to FRANK and RoSE}
\label{sec:transfer}

A metric calibrated on one benchmark may only have learned that benchmark: its systems, its annotators and its rating scale. We therefore tested whether the SummEval calibration of LNC transfers to two corpora with different systems and annotation schemes. The protocol below was fixed before any metric was run on them.

\subsection{Protocol}

LNC is applied with the frozen SummEval calibration, without refitting and without changing any feature. On FRANK its consistency scorer is correlated with the factuality score, and on RoSE its relevance scorer with the ACU score. The primary axis is the summary-level Spearman correlation with the copy rate partialled out; the raw correlation is secondary. The baselines are the cheap reference-free metrics NLIG, LGS, Lead-3 and Lead-5, together with UniEval and G-Eval for the matching dimension. G-Eval runs once with greedy decoding to fit the compute budget. The success criterion was stated as follows: LNC transfers if, on both FRANK CNN/DM and RoSE, (i) its partial correlation is significantly higher than that of the best cheap reference-free baseline in a paired bootstrap over articles, with Benjamini--Hochberg correction within the table of each corpus, and (ii) it is not significantly worse than G-Eval. FRANK XSum was declared an out-of-domain stress test outside the criterion, because its single-sentence summaries leave the order margin, LNC's most important perturbation feature, at zero.

AlignScore was added after the pre-registered run, as the strongest published consistency baseline we could run. It is reported in the tables but plays no role in the criterion. With it in the family of tests, the BH-adjusted $q$-values change slightly; none of the conclusions below depends on it.

One property of the criterion is worth stating before the results, because it governs how they should be read. NLIG is both a baseline in the family and the source of LNC's two entailment features. On a corpus annotated for consistency alone, the best cheap reference-free baseline in condition (i) is therefore LNC's own strongest component, and the condition asks not whether the frozen calibration transfers but whether the remaining eight features add significantly to it on a new corpus. That is a stricter question than the one the section sets out to answer. We did not foresee how severe the criterion would become on a single-dimension corpus, and we report it as written rather than rewriting it once the numbers were known.

\subsection{Results}

Tables~\ref{tab:transfer_frank_cnndm_confound}--\ref{tab:transfer_rose_cnndm_confound} give the raw and partial correlations on the three transfer sets, and Tables~\ref{tab:transfer_frank_cnndm_comparisons} and~\ref{tab:transfer_rose_cnndm_comparisons} the paired comparisons on the partial axis.

\paragraph{FRANK CNN/DM.} LNC reaches a partial correlation of .516 with the factuality score (raw .622), the highest of all metrics. It is significantly better than UniEval ($+.071$) and G-Eval ($+.070$), and far better than LGS and the lead controls. Its advantage over NLIG, the best cheap reference-free baseline, is $+.028$ with a 95\% interval of $[-.005, +.062]$ and is not significant ($q=.106$ in the pre-registered family, $q=.124$ with AlignScore added). It is also tied with AlignScore ($+.015$). Condition (ii) holds, but condition (i) fails.

\paragraph{FRANK XSum.} On the abstractive XSum summaries the copy rate carries almost no signal (raw correlation with factuality .063), so raw and partial correlations coincide. LNC reaches .280 and ties NLIG (.255), UniEval (.264), G-Eval (.264) and AlignScore (.293). It remains far ahead of LGS and the lead controls, which depend on copying.

\paragraph{RoSE.} On content coverage LNC transfers poorly. Its partial correlation of .148 ties Lead-3 (.151) and Lead-5 (.138) and is significantly lower than that of UniEval ($-.096$) and G-Eval ($-.089$). It is still significantly better than NLIG, AlignScore and LGS, whose correlation with the ACU score is close to zero once copying is controlled for. Both conditions fail.

\paragraph{Verdict.} The two calibrated scorers behave differently, and the useful reading of the test is per scorer rather than in aggregate.

The consistency scorer transfers. On FRANK CNN/DM it attains .516, the highest partial correlation of any metric in the pool, significantly above both UniEval and the local G-Eval judge, and level with AlignScore, a model trained specifically for factual alignment. On FRANK XSum, where the summaries are a single abstractive sentence and the order margin is identically zero, it still ties NLIG, UniEval, G-Eval and AlignScore, while LGS and the lead controls collapse. Condition (ii) holds on FRANK. Condition (i) fails there by $+.028$ over NLIG, $q=.106$ --- that is, the frozen calibration carries its accuracy to a new corpus, new systems and a new annotation scheme, but its margin over its own entailment component is not resolvable at this sample size. Most of the FRANK performance is indeed the NLI features, which reach .488 alone; the calibration adds to them without being separable from them.

The relevance scorer does not transfer. Its .148 on RoSE ties Lead-3 and Lead-5 and is significantly below UniEval and G-Eval, so both conditions fail. The reason is a mismatch of constructs rather than of corpora: SummEval relevance is judged against the source by annotators who read it, whereas the ACU score measures coverage of a reference. A reference-free metric calibrated on the former has no access to the latter, while UniEval reads the reference for relevance and G-Eval reads the full source with a much larger model. We take this as a bound on the metric's scope, not as a defect to be patched.

By the pre-registered criterion as a whole, which requires both conditions on both corpora, LNC does not transfer. We state that verdict in the terms we committed to, and we do not restate the criterion to fit the outcome. Its practical content is the two paragraphs above: the consistency scorer is usable on new news corpora, the relevance scorer is not usable outside the construct it was calibrated on, and the test was unable to separate the full composite from its entailment component on a consistency-only corpus. FRANK and RoSE have now served as test sets; any change to LNC prompted by these results would have to be validated on a further corpus, and a sharper comparison against NLIG would need a corpus with more than one dimension.
